\section{Problem Formulation}

Let $\bx_{i,t} \in \R^M$ denote the feature vector for asset $i$ at week
$t$, and let $r_{i,t+1}$ denote the realized return over the following
week.
Our goal is to learn a scoring function $f_\theta : \R^M \to \R$ such
that $f_\theta(\bx_{i,t}) > f_\theta(\bx_{j,t})$ whenever
$r_{i,t+1} > r_{j,t+1}$, for all $(i,j)$ pairs in the cross-section
at week $t$.
We then construct a long-short portfolio by going long the top decile
and short the bottom decile of $f_\theta(\bx_{\cdot,t})$ at each week.

\section{Rank-Normalized ListNet Loss}
\label{sec:method:loss}

Standard ListNet \citep{cao2007learning} converts scores and targets into
probability distributions via softmax with temperature $\tau$:
\begin{equation}
  P_i = \frac{\exp(s_i / \tau)}{\sum_j \exp(s_j / \tau)}, \quad
  Q_i = \frac{\exp(y_i / \tau)}{\sum_j \exp(y_j / \tau)},
  \label{eq:listnet_raw}
\end{equation}
where $s_i = f_\theta(\bx_{i,t})$ and $y_i = r_{i,t+1}$.
The loss is the cross-entropy $\cL = -\sum_i Q_i \log P_i$.

\paragraph{Problem with raw returns.}
When $y_i = r_{i,t+1}$, the target distribution $Q$ is dominated by a
few extreme-return assets in volatile weeks, and nearly uniform in calm
weeks.
This scale sensitivity destabilizes training across the heterogeneous
volatility regime of cryptocurrency returns.

\paragraph{Rank-percentile normalization.}
Following \citet{poh2021building}, we replace raw returns with
cross-sectional rank percentiles:
\begin{equation}
  \tilde{y}_{i,t} = \frac{\mathrm{rank}(r_{i,t+1};\, \{r_{j,t+1}\}_{j=1}^N) - 1}{N - 1}
  \in [0, 1],
  \label{eq:rank_norm}
\end{equation}
so the top-performing asset always receives $\tilde{y} = 1$ and the
worst-performing always receives $\tilde{y} = 0$, regardless of the
magnitude of the week's returns.
We substitute $\tilde{y}_{i,t}$ for $y_i$ in \cref{eq:listnet_raw}.
This ensures every asset contributes equally to the target distribution
and the loss gradient is invariant to cross-sectional volatility scaling.

\section{LSTM Architecture}
\label{sec:method:lstm}

The LSTM model encodes the $L$-step feature history for each asset
before performing cross-asset comparison.
Given the feature sequence $X_i = (\bx_{i,t-L+1}, \dots, \bx_{i,t}) \in \R^{L \times M}$,
the forward pass is:

\begin{enumerate}[nosep]
  \item \textbf{Input projection and position encoding.}
        $H = \mathrm{Dropout}(W_\text{in} X_i) + P$,
        where $W_\text{in} \in \R^{M \times d}$ and
        $P \in \R^{L \times d}$ is a learnable positional embedding ($d=64$).

  \item \textbf{LSTM encoder.}
        A single-layer LSTM processes $H$ to produce the final hidden
        state $h_i \in \R^d$.

  \item \textbf{Asset embedding.}
        A learnable asset embedding $e_i \in \R^d$ is added:
        $h_i \leftarrow h_i + e_i$.
        This allows the model to learn asset-specific biases.

  \item \textbf{Cross-asset attention.}
        All $N$ hidden states are stacked into
        $\mathbf{H} = [h_1, \dots, h_N]^\top \in \R^{N \times d}$.
        Multi-head self-attention (2 heads) produces
        $\mathbf{H}' = \mathrm{LayerNorm}(\mathbf{H} + \mathrm{MHA}(\mathbf{H}))$,
        enabling each asset to attend to all peers in the same week.

  \item \textbf{Output head.}
        $s_i = W_\text{out} h'_i \in \R$ is the ranking score for asset $i$.
\end{enumerate}

\noindent The model has approximately 53K parameters regardless of $M$,
as the number of features affects only the input projection layer.

\section{CrossSectionalGatedNet}
\label{sec:method:csgated}

The CrossSectionalGatedNet (CS-Gated) deliberately omits temporal
modeling.
At each week $t$, each asset's feature vector $\bx_{i,t}$ is processed
independently through a stack of $L = 3$ Gated Residual Network (GRN)
layers \citep{lim2021temporal}:
\begin{equation}
  \mathrm{GRN}(\bx) = \mathrm{LayerNorm}\bigl(\bx + \mathrm{GLU}(\mathrm{ELU}(W_1 \bx + b_1) W_2 + b_2)\bigr),
\end{equation}
followed by a cross-asset attention layer that allows each asset to
attend to all other assets in the same week.
The motivation is that, following \citet{gu2020empirical}, the cross-section
of contemporaneous features may already encode sufficient signal for ranking.

\section{Temporal Fusion Transformer}
\label{sec:method:tft}

The TFT follows \citet{lim2021temporal} with an encoder-only design
adapted for cross-sectional ranking.
The forward pass for asset $i$ proceeds as follows:

\begin{enumerate}[nosep]
  \item \textbf{Per-variable embedding.}
        Each of the $M$ input features is independently projected to
        $\R^d$ via a learned linear layer, yielding
        $\mathbf{V}_i \in \R^{L \times M \times d}$.
        A learnable positional embedding is added.

  \item \textbf{Variable Selection Network (VSN).}
        The corrected VSN processes each variable through a shared GRN
        (approximately 27K parameters), then computes softmax selection
        weights $\mathbf{w} \in \Delta^{M-1}$ conditioned on the asset
        embedding $e_i$ as static context.
        The selected representation is the weighted sum of GRN outputs.
        This design follows \citet{lim2021temporal} more closely than
        the commonly used lightweight linear VSN ($\approx$5K parameters).

  \item \textbf{LSTM encoder.}
        A single-layer LSTM produces temporal context from the
        VSN-selected features.
        A gated skip connection and LayerNorm are applied.

  \item \textbf{Interpretable multi-head attention.}
        TFT-style attention with a shared value projection and
        head-averaged outputs (2 heads) attends over the $L$
        LSTM outputs.

  \item \textbf{Cross-asset attention and output.}
        Identical to the LSTM: all $N$ asset representations are passed
        through cross-asset self-attention before a linear scoring head.
\end{enumerate}

\noindent The model has approximately 90--100K parameters, reduced from
209--368K in unconstrained TFT configurations through capacity regularization.

\section{Gradient Accumulation}
\label{sec:method:gradacc}

At each training step, the ListNet loss is computed over $N = 86$ assets
(one week).
We accumulate gradients over $K = 4$ consecutive weeks before updating
parameters, yielding an effective batch of $N \times K = 344$ asset-week
observations per update.
This stabilizes the rank estimate in \cref{eq:rank_norm}: with 86 assets,
the rank percentile has resolution $1/85 \approx 0.012$; pooling 4 weeks
improves rank diversity without introducing look-ahead bias (each week's
targets are computed independently).

\section{32-Seed Ensemble}
\label{sec:method:ensemble}

Each deep learning model is trained 32 times with independent random
weight initializations.
The final ranking score for each asset is the mean of the 32 predicted
scores, and portfolio construction proceeds on this mean score.

Ensemble averaging is critical for deep learning models in this setting
because per-seed Sharpe ratios are highly variable.
For LSTM on the \texttt{Price+Technical} configuration, individual seeds
produce Sharpe ratios with mean $+0.70$ and standard deviation $0.80$,
with a minimum of $-0.77$ and maximum of $+2.51$.
The 32-seed ensemble Sharpe ratio is $+1.45$—2.1$\times$ the per-seed
mean—demonstrating that diverse random initializations learn
complementary ranking patterns whose noise cancels in the average.
The most extreme case is LSTM on \texttt{+Trump+ETF}: the per-seed mean
is only $+0.26$, yet the ensemble achieves $+2.03$ (7.8$\times$
amplification).
This behavior conforms to the law of large numbers applied to the
prediction errors, which are partially uncorrelated across seeds.

TFT exhibits weaker ensemble gain because its larger capacity ($\approx$95K
parameters) causes seeds to converge to more similar local minima,
reducing the prediction diversity needed for noise cancellation.

\section{Traditional Baselines}
\label{sec:method:baselines}

Traditional models use only the current-period feature vector $\bx_{i,t}$
(no lookback window), consistent with their cross-sectional formulation.
At each test week, predicted scores are used for decile sorting.

\textbf{OLS.}
A pooled ordinary least-squares regression of next-week return on current
features, estimated on the training set.
No regularization; provides a linear, fully-interpretable baseline.

\textbf{ElasticNet.}
OLS with combined L1/L2 regularization.
The candidate grid is
\[
\alpha \in \{0.001, 0.01, 0.1, 1.0\}, \qquad
\ell_1\text{-ratio} \in \{0.1, 0.5, 0.9, 1.0\}.
\]
Selection is based on validation Sharpe ratio, so regularization is tied to
the portfolio-ranking objective rather than in-sample prediction error alone.

\textbf{PCA Regression.}
Features are projected onto their top $k$ principal components
($k \in \{1, 2, 3, 5\}$, chosen on validation), then OLS is applied
in the reduced space.

\textbf{PLS.}
Partial least-squares regression, which finds latent directions that
maximize covariance with the return target
($k \in \{1, 2, 3, 5\}$ components, chosen on validation).

\textbf{Random Forest (RF).}
An ensemble of 300 regression trees with maximum depth $\in \{1,2,3,5\}$
and feature fraction $\in \{\sqrt{M}, 0.33, 0.5\}$.
Best hyperparameters are selected on the validation set.
Thirty-two independent forests (different \texttt{random\_state}) are trained
and their predictions averaged.

\textbf{Gradient Boosting (GBT).}
Gradient-boosted regression trees with depth $\in \{1, 2\}$,
learning rate $\in \{0.01, 0.1\}$, and 100 or 300 estimators.
As with RF, predictions are averaged over 32 seeds.

Linear models (OLS, ElasticNet, PCA, PLS) are deterministic given their
hyperparameters and are equivalent to a 1-seed ensemble.
The asymmetry between DL (32 seeds), tree models (32 seeds), and linear
models (1 seed) reflects standard deployment practice for each class;
Section~\ref{sec:analysis:ensemble} discusses its implications.

\section{Evaluation Metrics}
\label{sec:method:eval}

At each test week $t \in \mathcal{T}_{\text{test}}$, we form a long
(top decile) and short (bottom decile) portfolio based on predicted scores.
Portfolio weights are either equal-weight (EW) or price-weight (PW).
The annualized Sharpe ratio is
$\SR = \bar{r}_\text{LS} \sqrt{52} / \hat{\sigma}_\text{LS}$,
where $\bar{r}_\text{LS}$ and $\hat{\sigma}_\text{LS}$ are the mean and
standard deviation of weekly long-short returns.
Statistical significance is assessed via Newey--West $t$-statistics
\citep{newey1987simple} with automatic lag selection \citep{newey1994automatic},
and 95\% Sharpe-ratio confidence intervals are obtained by circular block
bootstrap \citep{politis1994stationary}.
The bootstrap block length is $b=\max(2,\lfloor T^{1/3}\rfloor)$, where
$T$ is the number of valid weekly portfolio-return observations.
